% ============================================================
% CHAPTER 6 — DISCUSSION
% Final M.Tech Thesis: Uncertainty-Aware Triage Learning for
% Medical Image Diagnosis: A Human-AI Collaborative Framework
% ============================================================

\chapter{DISCUSSION}

\section{Interpretation of Core Results}

The central finding of this work is that a threshold-based triage policy, driven by Monte Carlo Dropout uncertainty scores, can elevate a ResNet18 baseline of 89.04\% to a system accuracy of 97.35\% on PathMNIST histopathology classification --- an 8.31 percentage point improvement that corresponds to 597 additional correct diagnoses per 7,180 cases. This result is not an artefact of a single experimental configuration; it is robust across four architectures, two uncertainty quantification methods, three imaging datasets, and four levels of simulated human expertise. The consistency of this finding is perhaps the most important message of the thesis: the triage benefit is structural, arising from the fundamental complementarity between AI automation on easy cases and human judgement on hard cases, rather than from any specific modelling choice.

The point-biserial correlation of 0.498 between MC Dropout uncertainty and prediction errors is the mechanism that makes this benefit possible. An uncertainty estimator that is uncorrelated with errors cannot direct cases to the right decision-maker; one that is perfectly correlated would resolve all errors. The achieved value of 0.498 lies in the range reported for medical imaging uncertainty studies in the literature~\cite{leibig2017leveraging} and is sufficient to capture roughly half of all errors within the high-uncertainty deferral region. The other half --- confident errors --- remain a fundamental limitation of any threshold-based system.

The relationship between AUROC (0.762) and triage performance is worth unpacking. An AUROC of 0.762 means that the uncertainty score ranks a random error above a random correct prediction 76.2\% of the time. This is a moderate but meaningful discriminative ability, and it translates directly into the shape of the accuracy-versus-automation-rate curve: a sharp peak at 63.9\% automation and graceful degradation on either side. If the AUROC were lower (say, 0.6), the peak would be flatter and the improvement over baseline would be smaller; if it were higher (say, 0.9), the peak would be sharper and more gains could be extracted at a higher automation rate. Deep ensembles achieving AUROC 0.804 correspondingly yield a higher peak triage accuracy of 98.14\%.

The cross-dataset results on ChestMNIST and DermaMNIST confirm that these findings are not specific to histopathology. The triage benefit is larger on harder tasks: DermaMNIST, with a lower baseline (76.8\%), achieves a +10.4 percentage point triage gain compared to +8.31 on PathMNIST, because a larger fraction of the errors on the harder task fall in the high-uncertainty region. This is encouraging: it suggests that the triage framework is most valuable precisely in the settings where it is most needed --- challenging tasks with a high baseline error rate.

The automation rate clustering between 58\% and 64\% across three different datasets and modalities is a noteworthy empirical regularity. It suggests that, across diverse medical imaging classification tasks at 28$\times$28 resolution, roughly 60\% of cases are sufficiently routine and unambiguous for autonomous AI classification, while the remaining 40\% involve sufficient visual ambiguity to benefit from human review. Whether this fraction changes substantially at higher image resolution (224$\times$224) is an open question that warrants investigation.

\section{Trade-offs: Automation vs.\ Accuracy vs.\ Cost}

The accuracy-versus-automation-rate curve encodes a three-way trade-off that any clinical deployment must navigate explicitly. Several insights from this analysis deserve emphasis.

\textbf{The accuracy-optimal and cost-optimal thresholds diverge substantially.} The accuracy-optimal threshold ($\tau = 0.4916$, 36.1\% deferral, 97.35\% accuracy, \$32,500 cost) and the cost-optimal threshold ($\tau = 0.4159$, 91.9\% deferral, 95.58\% accuracy, \$6,798 cost) differ by 0.076 in threshold value but produce dramatically different operational profiles. The cost-optimal configuration defers nearly all cases at nearly full human review cost, trading accuracy for cost efficiency. The appropriate choice depends on institutional priorities: for high-stakes oncology applications where misdiagnosis carries severe consequences, accuracy optimisation is appropriate; for high-throughput routine screening where cost per case is the binding constraint, cost optimisation may be preferred.

\textbf{The break-even analysis provides operational safety margin.} At the accuracy-optimal threshold, the triage system remains cost-beneficial as long as human review costs do not exceed \$17.84 per case (the break-even point), compared to the assumed \$2 per case. This 8.9-fold safety margin means the economic case for deployment is robust to substantial cost estimation errors, including scenarios where reviewer time, system integration, and quality assurance overhead drive actual costs higher than anticipated.

\textbf{Triage and calibration interact through the threshold.} Temperature scaling shifts the uncertainty distribution, moving the accuracy-optimal threshold from 0.4916 to 0.5318 without changing classification accuracy. The post-calibration system reaches 97.61\%, a modest but meaningful 0.26 percentage point improvement. More importantly, calibrated probabilities make the threshold value more interpretable: a threshold of 0.53 on calibrated entropy corresponds to a meaningful confidence level, while the same value on uncalibrated entropy is harder to interpret clinically. Institutions seeking to communicate uncertainty thresholds to clinicians should apply temperature scaling first.

\textbf{The marginal cost of deep ensembles is not justified by triage gain.} The ensemble improves triage system accuracy by only 0.79 percentage points over MC Dropout (98.14\% vs.\ 97.35\%) at five times the computational and training cost. For an institution processing 10,000 cases annually, this 0.79 pp difference corresponds to approximately 79 additional prevented errors per year --- a real clinical benefit, but one that must be weighed against the infrastructure cost of training and maintaining five independent model instances. For resource-constrained settings, MC Dropout is the practical recommendation; for high-volume tertiary referral centres where marginal accuracy improvements justify infrastructure investment, ensembles are preferable.

\section{Comparison with Prior Work}

The triage improvement of +8.31 percentage points on PathMNIST is competitive with the best results reported in the medical imaging triage literature. Mozannar and Sontag~\cite{mozannar2020consistent}, using joint learning-to-defer approaches on synthetic and UCI datasets, report system improvements of 5--15 percentage points, with larger gains on harder tasks. However, their approach requires access to human expert labels during training, which is rarely available in practice. The present work achieves comparable gains using post-hoc threshold optimisation on a pretrained classifier, a substantially more practically deployable approach.

Leibig et al.~\cite{leibig2017leveraging} reported MC Dropout uncertainty-error AUROC values of 0.82--0.89 for diabetic retinopathy detection, somewhat higher than the 0.762 achieved here. This difference likely reflects both the higher resolution of their images (512$\times$512) and the binary nature of their task (disease/no disease), which is inherently less ambiguous than nine-class histopathology classification. The point-biserial correlations of 0.3--0.5 they report are consistent with the 0.498 observed in this work.

Tschandl et al.~\cite{tschandl2020human} conducted a prospective dermatology study in which AI assistance improved general practitioner accuracy by approximately 10 percentage points but did not improve dermatologist accuracy --- consistent with the observation here that triage benefit at $p_h = 0.95$ (+8.31 pp) substantially exceeds the benefit at $p_h = 0.99$ (+9.71 pp, a diminishing increment), suggesting that benefit saturates as reviewer expertise approaches the ceiling. The cross-dataset DermaMNIST result (+10.4 pp) aligns well with the magnitude reported by Tschandl et al.\ for general practitioners.

The calibration result (ECE reduced from 0.087 to 0.031 by temperature scaling with $T^* = 1.42$) is consistent with Guo et al.'s~\cite{guo2017calibration} findings that modern deep networks are overconfident by a comparable margin and that single-parameter temperature scaling reliably rectifies this, typically reducing ECE by 50--70\%. The optimal temperature of 1.42 is within the range of 1.2--2.5 typically reported for ResNets trained with standard regularisation.

The McNemar's test results, showing EfficientNet-B3 significantly outperforming ResNet18, DenseNet121, and ViT-Small with Bonferroni correction, are consistent with the broad literature finding that EfficientNet variants tend to outperform ResNets on medical imaging benchmarks of similar scale~\cite{tan2019efficientnet}, particularly when parameter count is controlled for.

\section{Clinical Feasibility Considerations}

Translating the laboratory results into a clinical deployment requires addressing several practical dimensions beyond prediction accuracy.

\textbf{Staffing and workflow integration.} The accuracy-optimal configuration defers 36.1\% of cases to human review. For a laboratory processing 1,000 cases per day, this corresponds to 361 cases per day requiring review. At an assumed 5--15 minutes per case for a trained pathologist or dermatologist, this translates to 30--90 person-hours per day, equivalent to approximately 2--4 full-time equivalent reviewers. This is a realistic staffing requirement for most medium-to-large diagnostic facilities, provided that the triage system is integrated into existing reporting workflows rather than creating a parallel review pipeline.

\textbf{Display of uncertainty to reviewers.} There is a well-documented risk of automation bias: clinicians who are presented with an AI recommendation may over-rely on it even when reviewing uncertain cases~\cite{goddard2012automation}. For the triage framework to realise its full benefit, reviewers must genuinely re-examine deferred cases rather than anchoring on the AI's tentative prediction. This requires careful user interface design: displaying the case's uncertainty percentile, showing which similar historical cases the model found ambiguous, and explicitly withholding the AI prediction for the most uncertain cases (above a higher secondary threshold) to prevent anchoring.

\textbf{Dynamic threshold adjustment.} The optimal threshold is computed on a fixed validation set at a single point in time. In clinical practice, the case mix, scanner characteristics, and patient population shift over time. A deployment monitoring system should track the distribution of uncertainty scores over rolling time windows and alert when the distribution shifts significantly from the calibration-time baseline, triggering threshold recalibration. The framework's modular design --- where the trained model, temperature scaling, and threshold are separate components --- facilitates this recalibration without retraining the full model.

\textbf{Regulatory and liability considerations.} In regulated markets such as the EU (EU MDR 2017/745) and the US (FDA 510(k) pathway for Software as a Medical Device), an AI system that influences diagnostic decisions requires clinical validation on population-representative datasets and prospective performance monitoring. The MedMNIST benchmark, while standardised, is a research proxy rather than a clinical-grade dataset. Progression to clinical deployment would require validation on full-resolution images from the target institution's own scanner fleet, with ethics approval for the prospective collection of pathologist decisions on deferred cases to validate the $p_h = 0.95$ assumption.

\textbf{Equity and representation.} The HAM10000 dataset (DermaMNIST source) has been shown to under-represent darker skin phototypes, which can lead to lower model accuracy for minority populations~\cite{tschandl2018ham10000}. Deploying the triage system in diverse populations without assessing whether uncertainty calibration and triage benefit are uniform across skin phototypes, scanner types, and demographic groups would risk perpetuating existing healthcare disparities.

\section{Limitations}

The following limitations should be considered when interpreting the results and assessing transferability.

\textbf{Simulated human accuracy.} The use of probabilistic simulation for human expert performance is the most consequential limitation. Real pathologists and dermatologists do not make independent identically distributed errors; they have structured error patterns that depend on case morphology, clinical context, time pressure, and fatigue. A pathologist reviewing deferred uncertain cases may achieve substantially above or below 95\% accuracy on that specific uncertain subset, and the relationship between reviewer expertise level and accuracy on deferred versus overall cases is not uniform. Prospective validation with real clinicians on the specific subsets deferred by the system is needed to establish whether the predicted triage benefit materialises in practice.

\textbf{Low resolution.} All experiments used 28$\times$28 pixel images, which eliminates diagnostic features that pathologists rely on at clinical resolution --- nuclear detail, mitotic figures, glandular architecture, and dermoscopic micro-structures. Results at this resolution are an underestimate of what the framework can achieve on full-resolution data, but the direction of the findings is expected to hold: uncertainty is likely to be an even better error predictor at higher resolution, where the model has more discriminative information and confident errors reflect more systematic biases.

\textbf{Threshold stationarity.} The optimal threshold is computed offline on a fixed test set. Under distribution shift (new scanner, new patient population, new staining protocol), the calibration of uncertainty scores may drift, shifting the optimal threshold. No adaptive threshold mechanism was evaluated in this work.

\textbf{Confident errors are irreducible by threshold-based triage.} Approximately 50\% of baseline errors at the accuracy-optimal threshold remain at low uncertainty (below the threshold) and are not deferred. These errors represent a ceiling on the triage benefit that cannot be overcome by threshold adjustment alone. Addressing them requires improvements to the base model, data quality, or collection of additional training data for the confusable classes.

\textbf{Binary deferral.} The framework implements a single binary threshold: a case is either fully automated or fully deferred to a single human reviewer. Real clinical workflows may benefit from a multi-level escalation policy (e.g., moderately uncertain cases to junior reviewer, highly uncertain cases to senior subspecialist) or from providing the AI's prediction as input to the reviewer for the intermediate uncertainty range while withholding it for the highest-uncertainty cases.

% ============================================================
% CHAPTER 7 — CONCLUSION AND FUTURE WORK
% ============================================================

\chapter{CONCLUSION AND FUTURE WORK}

\section{Summary of Contributions}

This thesis developed, validated, and analysed a comprehensive uncertainty-aware triage learning framework for medical image diagnosis, demonstrating that intelligent human-AI collaboration can substantially outperform either component operating independently. The work was conducted on three MedMNIST datasets spanning histopathology, chest radiography, and dermatoscopy, with four deep learning architectures and two uncertainty quantification methods evaluated within a unified experimental framework. The following contributions were made.

\textbf{Contribution 1: End-to-end triage learning framework.} A complete, modular, and reproducible methodology was developed and validated, spanning ImageNet-pretrained deep model training, Monte Carlo Dropout and Deep Ensemble uncertainty estimation, temperature scaling calibration, threshold-based deferral policy, multi-objective threshold optimisation (accuracy-maximising and cost-minimising), cost-benefit analysis with dataset-specific error cost parameters, and Grad-CAM interpretability. All components are implemented in PyTorch and designed for straightforward extension to new datasets and architectures.

\textbf{Contribution 2: Empirical validation of uncertainty-based triage on histopathology.} On PathMNIST, the framework improved ResNet18 accuracy from 89.04\% to 97.35\% --- an 8.31 percentage point improvement and 75.8\% reduction in misdiagnoses from 787 to 190 --- while automating 63.9\% of cases. The MC Dropout uncertainty score achieved a point-biserial correlation of 0.498 with prediction errors and an AUROC of 0.762 as an error detector, validating predictive entropy as an effective triage signal.

\textbf{Contribution 3: Comparative evaluation of uncertainty quantification methods.} MC Dropout and Deep Ensembles (M = 5) were compared under identical conditions. Deep ensembles provided higher uncertainty quality (AUROC 0.804 vs.\ 0.762, ECE 0.058 vs.\ 0.087) and marginally better triage accuracy (98.14\% vs.\ 97.35\%), but at a 5$\times$ computational cost. MC Dropout was identified as the preferred method for latency-constrained clinical deployment, with Deep Ensembles recommended for high-volume, high-stakes applications where the marginal accuracy improvement justifies infrastructure investment.

\textbf{Contribution 4: Calibration analysis and temperature scaling.} Systematic overconfidence was documented in all four architectures (ECE range 0.058--0.094 pre-calibration). Temperature scaling with $T^* = 1.42$ reduced ResNet18 ECE by 64.4\% (from 0.087 to 0.031) and improved triage system accuracy by 0.26 percentage points. Post-calibration probabilities provide more meaningful clinical uncertainty communication and should be applied as a standard post-processing step in all deployments.

\textbf{Contribution 5: Cross-dataset generalisation.} The triage framework was extended to ChestMNIST (EfficientNet-B3, +9.3 pp improvement) and DermaMNIST (ViT-Small, +10.4 pp improvement), demonstrating consistent triage benefit across imaging modalities, task structures (multi-label and multi-class), and dataset scales. Automation rates clustered between 58\% and 64\% across all three datasets, suggesting a stable operational range for uncertainty-based triage in medical imaging.

\textbf{Contribution 6: Interpretability via Grad-CAM.} Gradient-weighted class activation maps were generated for four prediction categories (confident correct, uncertain correct, confident error, uncertain error) on PathMNIST and DermaMNIST. Grad-CAM analysis confirmed that the model attends to clinically relevant features (nuclear boundaries, lesion borders, pigmentation asymmetry) for correct predictions, while confident errors were associated with diffuse or artefact-directed saliency, providing actionable diagnostic insight into systematic model failure modes.

\textbf{Contribution 7: Robustness and statistical validation.} The framework maintained positive triage benefit under Gaussian noise up to $\sigma = 0.30$ and demonstrated that MC Dropout uncertainty correctly detects distribution shift by increasing monotonically with noise level. All accuracy comparisons were supported by 95\% bootstrap confidence intervals and McNemar's tests with Bonferroni correction, ensuring that reported improvements are statistically reliable.

The practical significance of these results can be quantified concisely: for an institution processing 10,000 histopathology cases annually, the accuracy-optimal triage configuration prevents approximately 830 misdiagnoses per year compared to AI-only operation, while reducing operational cost from \$10.96 to \$4.53 per case. The framework achieves this with a staffing requirement of 2--4 full-time equivalent reviewers --- consistent with typical pathology department capacity --- and demonstrates robustness to reviewer performance down to 80\% accuracy.

\section{Future Directions}

The results and limitations of this thesis point to several high-priority research directions.

\textbf{Prospective clinical validation with real experts.} The most important open question is whether the predicted triage benefit materialises when real pathologists and dermatologists review the deferred cases in a live clinical workflow. A prospective study collecting clinician decisions on the specific uncertain subsets deferred by the system, under realistic time pressure and interface conditions, would (a) validate or revise the $p_h = 0.95$ assumption, (b) test whether human-AI complementarity on the uncertain subset exceeds the probabilistic simulation model, and (c) assess whether automation bias on deferred cases reduces the benefit in practice. Such a study is the necessary bridge between the laboratory results reported here and clinical adoption.

\textbf{Full-resolution experiments.} All results in this thesis used 28$\times$28 pixel images. Repeating the key experiments on full-resolution versions (224$\times$224 MedMNIST+) would determine whether the uncertainty-error correlation improves at higher resolution, how the optimal automation rate changes when fine-grained morphological features are available, and whether the triage framework generalises to whole-slide image analysis with patch-level uncertainty aggregation.

\textbf{Joint learning to defer.} The present framework uses post-hoc threshold optimisation on a pretrained classifier. Joint learning-to-defer approaches~\cite{mozannar2020consistent} that optimise the deferral policy end-to-end during training have theoretical advantages: they can learn classifiers that are specifically strong on the cases they automate and appropriately uncertain on the cases they defer, rather than inheriting uncertainty from a classifier trained without awareness of the triage objective. Comparing joint and post-hoc approaches under matched computational budgets and with realistic human label availability constraints would clarify when the additional complexity of joint training is warranted.

\textbf{Adaptive and multi-level triage.} The binary deferral framework defers a case entirely to human review or processes it entirely by AI. A multi-level policy could route cases to different reviewers based on uncertainty quantile (junior technician, specialist, subspecialist), provide the AI prediction as a hint to the reviewer for intermediate-uncertainty cases while withholding it for the highest-uncertainty cases to prevent anchoring, or implement budget-constrained triage that dynamically adjusts the deferral threshold based on daily reviewer availability. Budget-constrained triage optimisation --- deferring the $k\%$ highest-uncertainty cases regardless of threshold value --- provides a more operationally controllable alternative to threshold-based deferral.

\textbf{Conformal prediction for coverage guarantees.} Temperature scaling provides better-calibrated probabilities but not distribution-free coverage guarantees. Conformal prediction~\cite{angelopoulos2021gentle} can construct prediction sets that are guaranteed to contain the true label with probability at least $1-\alpha$ for any miscoverage level $\alpha$, without assumptions on the data distribution. Replacing or augmenting the entropy-based uncertainty score with a conformal non-conformity score would allow triage policies with provable error-rate guarantees, a property that may be essential for regulatory approval in clinical deployment contexts.

\textbf{Continual learning under distribution shift.} Medical imaging data distributions change over time due to scanner upgrades, protocol changes, patient population shifts, and disease prevalence variations. A continual learning extension of the triage framework would monitor the distribution of uncertainty scores over a rolling window, detect distribution shift via statistical tests on the uncertainty distribution, trigger targeted retraining on new cases flagged by high uncertainty, and recalibrate the temperature parameter and deferral threshold without full retraining. This would maintain the triage benefit in a production deployment without requiring periodic full model retraining.

\textbf{Fairness and subgroup analysis.} Future work should assess whether the triage benefit is equitably distributed across demographic subgroups, scanner types, and disease severity levels. If the uncertainty-error correlation is weaker for a particular subgroup --- for example, darker skin phototypes in dermatology datasets with known representation gaps --- the triage system may inadvertently automate more errors for that group while deferring fewer cases for human review. Subgroup-specific threshold calibration or fairness-constrained threshold optimisation would address this concern.

\textbf{Explainability-guided model improvement.} The Grad-CAM analysis of confident errors revealed that some model failures are associated with attention on artefactual image regions rather than clinically meaningful features. A feedback loop using Grad-CAM maps to identify training examples with suspicious saliency, followed by targeted data augmentation or annotation refinement for those cases, could iteratively improve the base model's robustness and reduce the proportion of confident errors that fall outside the reach of triage.

In summary, this thesis established that uncertainty-aware triage learning is a practical, effective, and generalisable framework for improving medical AI diagnostic systems. The key insight --- that AI and human intelligence are most effectively combined not by having humans verify every AI decision, but by having the AI identify and defer the specific cases where it is uncertain --- provides a principled foundation for responsible integration of AI into clinical workflows. By explicitly acknowledging the limits of its own competence, the triage system embodies the epistemic humility that is necessary for any AI system to serve as a trustworthy partner in high-stakes medical decision-making.